\documentclass[11pt, letterpaper]{article}

\usepackage[utf8]{inputenc}
\usepackage[T1]{fontenc}
\usepackage[left=1.7cm, right=1.7cm, top=1.7cm, bottom=1.7cm]{geometry}
\usepackage{xcolor}
\usepackage{hyperref}
\usepackage{titlesec}
\usepackage{enumitem}
\usepackage{tabularx}
\usepackage{charter}

% --- Color Palette ---
\definecolor{primary}{HTML}{003366}
\definecolor{secondary}{HTML}{444444}
\definecolor{rulecol}{HTML}{CCCCCC}
\definecolor{RULECOL}{HTML}{CCCCCC}

% --- Hyperlink Setup ---
\hypersetup{
    colorlinks=true,
    urlcolor=primary,
    linkcolor=primary
}

% --- Section Formatting ---
\titleformat{\section}
  {\large\bfseries\color{primary}\MakeUppercase}
  {}
  {0em}
  {}
  [\vspace{-6pt}\textcolor{rulecol}{\hrule height 0.8pt}\vspace{4pt}]

\titlespacing*{\section}{0pt}{14pt}{8pt}

% --- Custom Commands ---
\newcommand{\cventry}[4]{
    \noindent\textbf{#1} \hfill \textit{#2} \\
    \textit{#3} \hfill \textit{\small #4} \\
    \vspace{-4pt}
}

\begin{document}
\pagestyle{empty}

% === HEADER ===
\begin{center}
    {\Huge \textbf{Juan Pablo Arango Rentería}} \\
    \vspace{6pt}
    {\color{secondary} \itshape Research Analyst \quad|\quad Communication Engineer \quad|\quad Computational Neuroscience} \\
    \vspace{6pt}
    {\small
    Toronto, ON \quad\textbar\quad
    (613) 762-2123 \quad\textbar\quad
    \href{mailto:juan.pablo.arango.renteria@gmail.com}{juan.pablo.arango.renteria@gmail.com} \\
    \href{https://github.com/JPArangoRenteria}{github.com/JPArangoRenteria} \quad\textbar\quad
    \href{https://linkedin.com/in/juan-pablo-arango-0380b0191}{linkedin.com/in/juan-pablo-arango-0380b0191}}
\end{center}

\vspace{8pt}

% === RESEARCH INTERESTS ===
\section{Research Interests}

\begin{tabularx}{\linewidth}{@{}p{0.5\linewidth} X@{}}
    \textbf{Computational Neuroimaging:} MEG source connectivity, Diffusion Tensor Imaging (DTI), tractography, and structural parcellation. &
    \textbf{ML \& Clinical Translation:} Machine learning biomarkers for pain laterality, subtype characterization, and surgical outcome prediction. \\

    \textbf{Mathematical Modeling:} Graph Neural Networks (GNN), Explainable AI (XAI), stochastic modeling, and statistical learning frameworks. &
    \textbf{Digital Signal Processing:} Multi-modal sensor fusion, physiological time-series analysis, and clinical NLP for pain narratives.
\end{tabularx}

% === RESEARCH EXPERIENCE ===
\section{Research Experience}

\cventry{Hodaie Lab (PI: Dr. Mojgan Hodaie) --- University Health Network / University of Toronto}{Toronto, ON}{Research Analyst I}{Oct 2025 -- Dec 2026}

\begin{itemize}[leftmargin=1.5em, topsep=2pt, itemsep=2pt]
    \item Spearheading computational development of the lab's \textbf{diffusion neuroimaging}, and \textbf{structural radiomics} pipelines for biomarker discovery in trigeminal neuralgia (TN) and related neuropathic pain disorders.

    \item Developing \textbf{normative and predictive models} linking segment-wise diffusion microstructure features to clinical phenotypes including pain laterality and surgical outcomes.

    \item Assisting in the transition from a \textbf{DataLad-based neuroimaging data management framework} to a centralized \textbf{TrueNAS storage infrastructure}, supporting scalable imaging data organization and computational workflow integration.

    \item Collaborating across projects on tractography refinement, spatial parcellation architecture, and explainable AI visualization frameworks for chronic pain neuroimaging.
\end{itemize}

\vspace{6pt}

\cventry{GNN-Based Wireless Network Optimization --- Jie Gao Research Group (PI: Dr. Jie Gao), Carleton University}{Ottawa, ON}{Research Collaborator}{Apr 2025 -- Present}

\begin{itemize}[leftmargin=1.5em, topsep=2pt, itemsep=2pt]
    \item Designing algorithms combining \textbf{Graph Neural Networks (GNN)} and \textbf{Reinforcement Learning (RL)} for distributed radio resource scheduling in wireless networks.

    \item Developed Python simulation frameworks evaluating RF system performance under varying network conditions and traffic loads.

    \item Applied machine learning optimization techniques directly analogous to graph-based connectivity modelling used in computational neuroscience.
\end{itemize}

\vspace{6pt}

\cventry{Biomedical ML Pipeline --- CUBIC Lab (PI: Dr. Rafik Goubran), Carleton University}{Ottawa, ON}{Research Collaborator}{Apr 2025 -- Present}

\begin{itemize}[leftmargin=1.5em, topsep=2pt, itemsep=2pt]
    \item Applied machine learning methods to physiological signal datasets for automated clinical event classification downstream of DSP preprocessing pipelines.

    \item Developed feature engineering and statistical evaluation workflows for multi-modal biomedical datasets.
\end{itemize}

% === SELECTED RESEARCH PROJECTS ===
\section{Selected Research Projects (Work in Progress)}

\vspace{4pt}
\cventry{Trigeminal Neuralgia Tractometry \& Microstructural Biomarker Mapping}{Hodaie Lab}{Computational Analyst}{2025 -- Present}

\begin{itemize}[leftmargin=1.5em, topsep=2pt, itemsep=2pt]
    \item \textbf{Tractometry Pipeline Engineering:} Automated end-to-end diffusion tractometry workflows across a 120-patient cohort, executing along-tract sampling of trigeminal nerve bundles to capture localized microstructural variations beyond global nerve averages.

    \item \textbf{Along-Tract Feature Extraction:} Quantified segment-wise microstructural metrics—including \textbf{FA}, \textbf{MD}, \textbf{AD}, and \textbf{RD}—to map localized demyelination patterns, axonal integrity, and structural degradation along the nerve profile.

    \item \textbf{Statistical \& Phenotypic Modeling:} Modeled high-dimensional tractometric feature profiles against clinical phenotypes, identifying significant focal along-tract alterations associated with pain laterality and disease severity.

    \item \textbf{High-Dimensional Visualization:} Standardized spatial mapping of tract-profile statistics into 3D visualization environments (MRtrix3, 3D Slicer), translating quantitative tractometric arrays into intuitive pathological heatmaps for neurosurgical analysis.
\end{itemize}
\vspace{4pt}

\cventry{Peripheral Node Parcellation \& Clinical Mapping}{Collaborator: Sarasa Tohyamaa --- Hodaie Lab}{Computational Analyst}{November 2025 -- October 2026}

\begin{itemize}[leftmargin=1.5em, topsep=2pt, itemsep=2pt]
    \item Extending trigeminal spatial parcellation and heatmap frameworks to peripheral nerve nodes in chronic pain patients.

    \item Performing morphological analyses to identify structural vulnerability patterns associated with chronification and nerve damage.
\end{itemize}


\vspace{4pt}

\cventry{Biomedical DSP: Automated Apnea \& Hypopnea Detection}{CUBIC Lab}{ML Analyst}{Collaborators: Matthew Stewart \& Julien Larivière-Chartier}{2025 -- Present}

\begin{itemize}[leftmargin=1.5em, topsep=2pt, itemsep=2pt]
    \item Applied machine learning classification and statistical event-scoring methods to physiological signal features for automated apnea and hypopnea detection.

    \item Developed feature engineering workflows using spectral and amplitude-derived characteristics to identify respiratory desaturation events.
\end{itemize}

\vspace{4pt}

\cventry{Glossopharyngeal Neuralgia (GPN): A Multimodal Analysis Case Series }{Hodaie Lab}{Collaborators: Emilio Garcia Flores, Jana Mahmoud, Annabella Helvadjian, Jason Yue(Lead)}{Computational Neuroimaging Analyst}{January 2026 -- August 2026}

\begin{itemize}[leftmargin=1.5em, topsep=2pt, itemsep=2pt]

    \item \textbf{Cohort Design:} Constructed a matched-control framework using \textbf{4 healthy controls per patient} (GPN: $n=7$, HC: $n=28$) for statistically robust analysis in a rare neuropathic pain cohort.

    \item \textbf{Tractography Pipeline:} Developed diffusion MRI tractography workflows to isolate the \textbf{glossopharyngeal nerve (CN IX)}, addressing anatomical challenges due to its small calibre and proximity to adjacent cranial nerves.

    \item \textbf{Microstructural Analysis:} Performed segment-wise and voxel-wise analysis of diffusion metrics (FA, MD, AD, RD) to characterize structural deviations between GPN patients and controls.

    \item \textbf{Framework Extension:} Generalized trigeminal radiomics pipelines to additional cranial nerve systems for broader application to rare cranial neuropathies.
\end{itemize}

\vspace{4pt}



% === PUBLICATIONS & CONFERENCE CONTRIBUTIONS ===
\section{Manuscripts}

\begin{itemize}[leftmargin=1.5em, topsep=2pt, itemsep=4pt]
    \item Garcia Flores E, Wu M, Cabrera-Montes J, \textbf{Arango JP}, Jörgens D, Li J, Latypov TH, Stewart M, Srisaikaew P, Hodaie M. ``Ventral Tegmental Area resting state functional connectivity is altered in Trigeminal Neuralgia Patients.'' \textit{Brain}. In review.
    \item Cheng J, Zuo K, Naraghi A, Tohyama S, \textbf{Arango JP}, Baltzer H. ``Median nerve regeneration after acute versus delayed repair: longitudinal tractography and tractometry in a case of bilateral transection.'' \textit{Muscle \& Nerve}, In review.
\end{itemize}

\section{Conference Presentations \& Abstracts}

\begin{itemize}[leftmargin=1.5em, topsep=2pt, itemsep=4pt]
    \item Wu M, Garcia Flores E, Cabrera-Montes J, \textbf{Arango JP}, Jörgens D, Li J, Latypov TH, Stewart M, Srisaikaew P, Hodaie M. ``Ventral Tegmental Area resting state functional connectivity is altered in Trigeminal Neuralgia Patients.'' Oral presentation and abstract, World Society for Stereotactic and Functional Neurosurgery (WSSFN) Congress, Marseille, France, 2026.
    \item Yuen J, Garcia Flores E, \textbf{Arango JP}, Mahmood J, Tohyama S, Helvadijan A, Hodaie M. ``Quantitative Multimodal MRI Characterization of Glossopharyngeal Neuralgia: An Imaging-Enriched Case Series.'' Abstract, World Society for Stereotactic and Functional Neurosurgery (WSSFN) Congress, Marseille, France, 2026.
    \item Cabrera J, Villamil F, Jörgens D, \textbf{Arango JP}, Hodaie M. ``Trigeminal Neuralgia Tractometry: Insights from a 120-Patient Cohort.'' Presentation, World Society for Stereotactic and Functional Neurosurgery (WSSFN) Congress, Marseille, France, 2026.
\end{itemize}

% === EDUCATION ===
\section{Education}

\cventry{Carleton University}{Ottawa, ON}{B.Eng. Communication Engineering (Co-operative)}{Aug 2020 -- Jul 2025}

\begin{itemize}[leftmargin=1.5em, topsep=2pt, itemsep=2pt]
    \item \textbf{Distinctions:} Dean's Honour List (2021), George Fierheller Scholarship (2021).

    \item \textbf{Relevant Coursework:} Digital Signal Processing (A+), Introduction to Machine Learning (A-), Wireless Communications (A-), Topics in Communication Networks (A), Digital Communications (A-), Network and Software Systems (A), Signals and Systems (B+).

    \item \textbf{Capstone:} LEO Satellite Network Simulator --- RF link budget modeling, OSPF routing, and ML-based performance prediction for Low Earth Orbit systems.
\end{itemize}

\vspace{6pt}

\cventry{University of Waterloo}{Waterloo, ON}{B.Math. Applied Mathematics, Honours}{Aug 2016 -- Aug 2020}

\begin{itemize}[leftmargin=1.5em, topsep=2pt, itemsep=2pt]
    \item \textbf{Distinctions:} University of Waterloo Merit Scholarship (2017).

    \item \textbf{Program Focus:} Scientific computation, numerical methods, stochastic modeling, and applied mathematical analysis.

    \item \textbf{Selected Coursework:} Control Theory (84\%), Computational Mathematics (88\%), Computational Linear Algebra (81\%), Partial Differential Equations I \& II, Applied Real Analysis, Sampling Theory.
\end{itemize}

% === INDUSTRY EXPERIENCE ===
\section{Industry Experience}

\cventry{Cloud RAN Systems Designer --- Ericsson Inc.}{Ottawa, ON}{Intership}{May 2023 -- Aug 2024}

\begin{itemize}[leftmargin=1.5em, topsep=2pt, itemsep=0pt]
    \item Developed 5G RAN radio hardware/software systems involving TRX architecture, RF link budgets, Opensearch monitoring infrastructure, and CI/CD pipelines using Docker, Kubernetes, and Jenkins.
\end{itemize}

\vspace{4pt}

\cventry{Embedded Software Engineer Intern --- Nokia}{Ottawa, ON}{Internship}{May 2022 -- Aug 2022}

\begin{itemize}[leftmargin=1.5em, topsep=2pt, itemsep=0pt]
    \item Optimized C/C++ real-time communication protocols achieving 25\% latency reduction and developed network performance testing tools in Unix environments.
\end{itemize}

\vspace{4pt}

\cventry{Full Stack Developer --- Neurobotx}{Ottawa, ON}{Contract}{Aug 2021 -- Dec 2021}

\begin{itemize}[leftmargin=1.5em, topsep=2pt, itemsep=0pt]
    \item Developed real-time biomedical device communication systems and REST API interfaces for data acquisition and control.
\end{itemize}

% === TECHNICAL SKILLS ===
\section{Technical Skills}

\begin{description}[style=multiline, leftmargin=3.5cm, font=\bfseries, itemsep=3pt]

    \item[Neuroimaging] MNE-Python, FreeSurfer, FSL (dtifit, eddy), MRtrix3, 3D Slicer, PyRadiomics, SimpleITK, FSLEyes, MIBrain, tractography, diffusion MRI processing.

 
    \item[Signal Processing] Spectral analysis, STFT, digital filter design (SciPy), physiological signal analysis, multi-modal sensor fusion, DSP pipelines.

    \item[Programming] Python, MATLAB, R, C/C++, JavaScript, Bash/Shell, SQL, Git, Linux/Unix.

    \item[Infrastructure \& DevOps] Docker, Kubernetes, Jenkins, AWS, CI/CD, Prometheus, Grafana, Opensearch, DataLad, TrueNAS.

\end{description}

% === LANGUAGES ===
\section{Languages}

\noindent
\textbf{Spanish} (Native)
\quad\textbar\quad
\textbf{English} (Fluent)
\quad\textbar\quad
\textbf{French} (Conversational)
\hfill
